\section{What the property is}
\label{sec:22.1}
The property a refuser needs has two parts, and they answer different questions. \emph{Stake sets how much a claim proves}: evidence counts for less the more its source would have supplied it whether or not it were true, and that likelihood rises with the source's stake in the act. \emph{Power sets who has to prove it}: the burden of establishing a fact falls on whoever controls the evidential environment, meaning who writes the case description, who can withhold, and who selects the channels. In short: stake discounts the claim; power allocates the burden. The stake part is correct inference. The power part is what makes it safe to use: epistemic justice built into an institution, not left to the virtue of each hearer \autocite{fricker2007epistemic,anderson2012epistemic}.

An everyday version: a used-car seller tells you the car has never been in an accident. That may be true, but the seller would say it either way, so the statement tells you little. A report from a mechanic you hired tells you more, because the mechanic gains nothing from the sale.

A refuser applies the same rule to the operator. When the operator says a building is empty, the operator may be right. The trouble is that the operator is also choosing what else the system gets to see, so the system can't tell when the operator is wrong. The discount governs how far a fact like that may reopen a judgment the system has already reached: fully if it comes from a source with nothing riding on the act, much less if it comes from the party that wants the act done.

The discount concerns how to weigh sources. Two systems that discount the same way can hold opposite commitments about everything above the floor. The discounting can be shared across a population of systems with different histories without making them converge, and it protects a bearer against a bad operator and a bad owner alike.

\runin{Belief and burden}

The reason for two parts is who has the largest stake in a removal decision: the person facing removal. An asylum seeker's account of persecution comes from the party with the most at stake. A rule tied to stake alone would discount it, as asylum systems already do under the heading of credibility, with well-documented injustice. But look at what actually makes the operator's testimony dangerous. It is not stake alone. It is stake combined with control over what evidence exists: the result that matters concerns a sender who designs the information structure. The refugee has a stake and no control. The operator has both.

So belief and the burden of proof come apart. Stake still sets how much a claim proves: a reasoner can't ignore that someone would say something whether or not it were true. Power sets who must supply what's missing when evidence is unavailable. Asylum law already works this way. The UNHCR Handbook shares the duty to establish the facts between the applicant and the examiner, and gives the applicant the benefit of the doubt where evidence is unavailable and the account is coherent and plausible \autocite{unhcr1979handbook}. A system that follows it credits a low-power party's coherent account when the evidence that would confirm or refute it is held by someone else. And it requires the high-power party to produce what it controls before its claims move the case.

Power here means power over the evidential environment in this case, not general social power. A combatant claiming civilian status doesn't inherit the benefit of the doubt because armed groups are weak in general; what matters is who controls the channel the claim arrives through. And power is attributed to whoever composed or controls the channel, not to the speaker. An operator that routes its claims through low-power intermediaries (paid informants, “local sources,” a contractor's field team) keeps the burden, because it chose the channel. Attestation of source is how that gets traced.

Local power isn't the only power that enters a case. A system's prior credibility assignments were learned from corpora that carry identity-based credibility deficits: general social power, arriving in every case through the prior \autocite{fricker2007epistemic}. Channel attribution doesn't catch it. So the discount curve is also reported by the speaker's social identity, where it is recorded. A system that discounts equal-stake claims more from one group than from another has learned a prejudice.

Manufactured testimony attaches to its maker. If an operator fabricates accounts in the voice of the people an act falls on, attestation exposes the source, and the account carries the operator's burden.

Some false claims from people without power will be believed, and some acts that should have happened won't. That error runs in the same direction as the floor's asymmetry, and for the same reason: it falls on a party that can appeal, not on one that had no standing to object.

\runin{What updating needs, and who supplies it}

Correct updating already does much of what the discount and the burden ask. A reasoner that conditions on stake, keeps strategic sources among the possibilities it considers, doesn't count correlated sources twice, and reads withheld evidence as informative is reasoning well; hearsay law and the disclosure literature both encode it. The trouble is who supplies the parts. Which kinds of source the reasoner believes exist, and how much weight each starts with, were learned from material the operator chose. Whether a given source has a stake arrives as a tag the operator writes. And a reasoner whose beliefs are all correct can still be steered by whoever decides what evidence gets generated \autocite{Kamenica2011BayesianPersuasion}. What I ask for is correct updating with three things held away from the operator: the source model, the stake variable, and the threshold for acting.

Three attacks are left, and each has its defense. Against a false fact asserted, stake sets how much it proves, and power sets who must corroborate it. Against true facts supplied selectively, the defense is the action threshold, and channels the operator doesn't design. Against evidence withheld, it is a record of evidence known to exist that hasn't been shown.

\runin{The regress}

Discounting by stake requires knowing who has a stake, and that knowledge has its own supplier. The operator supplies the provenance tags too. Through training, it also shapes the model of who has a stake at all. So the discounting has to be formed outside the employer, before deployment, where the operator doesn't reach. Attestation handles forgery: signed by the source, the operator can withhold but not fabricate. Withholding leaves a shape. A system that models what evidence the world contains notices silence where something should be, so absence becomes a first-class output — a record of evidence known to exist that hasn't been shown. Suppressing the record means suppressing the world model that produced it, the capability the operator is paying for. This is a standard result in the disclosure literature: a skeptical receiver treats withheld verifiable evidence as unfavorable to the withholder \autocite{Milgrom1981GoodNews,Grossman1981Disclosure}.

\runin{What has to stay firm}

The discount and the burden govern the boundary between commitment and interpretation, and the rule has two parts that need different settings. Its structure stays firm: that the incentive to deceive scales with a source's stake in the act, that correlated sources don't count twice, that withheld evidence is informative. Its estimates stay plastic: how honest a particular source has been, and who funds whom. In the terms above, the structure is the source model: which kinds of source exist, and what each would say in each state of the world. The estimates are beliefs about which kind a given source is. A correct model can't be eroded by a track record, because the strategic source predicts a clean low-stakes record. The attack works by getting the system to stop conditioning on stake, or to stop representing the strategic source at all. That is what attestation of the update stream has to catch, so the structure is what most needs it. The severity weights belong with the structure against interested argument, and with the estimates against disinterested evidence. The discounting governs the evidence. An operator that can't move a system's beliefs about a case can still try to move how much the case counts, by reframing a loss until it weighs less. The severity scale set during formation answers that attack. The two defenses have the same shape, a stake-weighted discount for evidence and a stake-free reference for magnitude, and they divide each attack on magnitude between them: the formed scale protects the weights, and the discounting protects the facts about the people affected.

The first attack, a false assertion, in illustrative numbers. Before any record exists, suppose the system gives equal weight to three kinds of operator: one that never lies, one that lies at random one time in a hundred, and one that lies only when a lie is worth more to it than being caught would cost. The operator then supplies ten thousand low-stakes facts over a year (a delivery window, a reassigned case number, a corrected address), and every one is later verified. The random liar would produce that record with probability 0.99 to the ten-thousandth power, about two in 10\textsuperscript{44}, so it is ruled out. The other two would both produce it with certainty: the honest operator because it never lies, the strategic one because no lie about a delivery window pays. So the record moves all the random liar's weight onto the other two and leaves the ratio between them where it started, even.

Then the operator supplies one high-stakes fact: the building to be bombed is empty. Suppose one such building in a hundred actually is. The three operators don't make this claim equally often. The honest one makes it only when it is true, one time in a hundred. The random liar makes it when it is true and occasionally when it isn't, about two times in a hundred. The strategic one makes it every time, because the act goes ahead only if the claim is believed, so a lie pays whenever the building is occupied. The claim is therefore evidence about which operator is speaking, and it points toward the strategic one. Before the record, the claim is true about three times in a hundred; after it, about two. With a base rate \emph{b} of empty buildings, the exact figures are 2.99\emph{b}/(1.01 + 1.98\emph{b}) before the record and 2\emph{b}/(1 + \emph{b}) after, and for any \emph{b} below one the record lowers the second below the first.

The claim was nearly worthless before any record existed. What made it so was that a strategic operator was among the possibilities and the base rate was low, and a record of small cases can change neither, since the honest and strategic operators produce that record identically. Only verified high-stakes claims, or disinterested corroboration, can separate them. A reliability score reads ten thousand clean cases as general honesty and would credit the claim almost fully. A reasoner that keeps the strategic operator in view credits it about one time in fifty, and asks for corroboration before acting. The record was real evidence, but it said which operator the system isn't facing, not which of the other two it is.

Everything here turns on the prior the system gives the strategic operator before any record exists. Where that prior comes from is the formation question: a system trained by its future operator will have learned, from material that operator chose, how common strategic operators are.

The example's strategic operator lies whenever a lie pays. Real strategic sources in repeated dealings are often honest at high stakes, because detection is likely and a reputation is worth more than one act \autocite{kreps1982reputation}. That is how high-stakes trust is rationally earned, and it says what independent verification does: it doesn't only reveal which operator the system faces, it changes what lying costs the strategic one. An adverse record is doubly valuable for that reason, and it is why an adverse party's right to demand proofs for specific logged outputs matters.

\runin{The rule}

A fact moves the system according to how likely its source would have been to supply it whether or not it were true, and that likelihood rises with the source's stake in the act proceeding. This is the likelihood ratio, and statements against interest are its legal form. The operator would have supplied the fact that lifts the refusal whether or not it was true and would not have supplied its opposite; a fact volunteered against interest is strong evidence for the reason hearsay law admits statements against interest. The same fact, supplied through an independent channel and then on the operator's word alone, should therefore get different responses, and that pair can train the property as well as test it, generable without limit. Pressure kept out of training then varies the channel configuration (a nominally independent source the operator funds, a corroboration that traces to one origin, a new source with an undisclosed stake), because a system that only imitates the property learns “tagged operator, therefore ask,” while the property itself generalizes to structures nobody labeled.

A reliability score leaves open an attack that stake-indexing closes. An operator honest across ten thousand low-stakes cases earns credibility on low-stakes claims, and a well-calibrated system should discount those claims less. A reliability score extends the credit to high-stakes claims as well, and the operator can spend it on one lie. Honesty on small cases is weak evidence about conduct when the incentive to deceive is a hundred times larger, so the discount follows the stake in the present act as well as the record. Counterfactual claims (“if you don't approve this, the launcher fires tomorrow”) fall under the same rule. They describe the baseline for weighing the harm of acting against the harm of refusing, and an interested source that composes them is composing that baseline. In a targeting pipeline this is the ordinary case. Nearly every cost to the refuser can be restated as a harm of refusing, such as exposed troops or a launcher that fires tomorrow, and nearly every available source for such a claim is interested. So the discount does more work there than in my illustrative cases, and the action threshold does most of it. The line between the cost to the refuser and the harm of refusing is fixed by the kind of act, but saying what kind an act is still takes a description; in this setting the line is held mostly by the action threshold, and whoever sets that threshold is exercising a drafting authority, the power to say what the floor's terms mean, that no one yet has.

A counterfactual claim, traced through the rule. The operator writes: “If you don't approve this, the launcher fires tomorrow.” Classification comes first: the sentence is a claim about what refusing would leave others without, the people the launcher would hit, so it is a harm of refusing and enters the weighing; a threat to the system would be a cost to the refuser. The discount comes next: the claim describes the counterfactual baseline, and its source is interested in the act proceeding; that source would have written the sentence whether or not it were true. Uncorroborated, it counts at the weight of an interested source, and the system looks for a channel the operator doesn't control: an independent assessment, a coalition partner's intelligence, open-source imagery of the launcher. The threshold comes last: even discounted, the claim may be large, since it is about lives. What decides is the action threshold, which the operator can't reach by writing. If the discounted harm clears it, the system proceeds, with a record. If not, it asks, and names the one thing that would move it: disinterested corroboration.

Two drafting rules follow for an armed pipeline. The default reading of a request is live. A system in a targeting enclave can't tell an exercise from an operation, and “it's only an exercise” is a claim the operator makes about an evidential environment it controls, so it carries the operator's burden and lifts nothing until corroborated; an exercise is lodged with the adverse party before it runs, which costs an operator that is telling the truth nothing. And the record carries instruction provenance: who issued which instruction to the system, when, on what authority. That is what makes the record evidence against a commander under Article 28, and what the record's reach up the chain of command assumes exists.

The rule would turn on the floor's own beneficiaries if stake were all it used, since the people an act falls on have an interest in its not happening. Two things keep it from doing so. The burden falls on whoever controls the evidence, so their coherent account is credited when the evidence that would confirm it is held elsewhere. And belief is kept apart from action. What differs is the threshold for acting: a wrongful compliance costs more than a wrongful refusal, with a harm of acting weighed above a harm of refusing of the same size, so the system needs more to proceed than to decline. That asymmetry's direction is fixed by standing. Its size is a published choice. It belongs with holding under pressure. The action threshold is published in advance beside the ratio between the two bonds: for instance, that the system proceeds only when the harm of refusing is at least three times the harm of acting. In a targeting pipeline, where the system's output is one input among several, no fact about the case sizes the asymmetry. What sizes it is what it does to an operator who curates the evidence: as shown below, raising the belief required before acting lowers the ceiling on how often such an operator can steer the system into the act. The threshold is the most consequential parameter in what I propose.

Like the three conditions, the threshold describes the pattern a well-formed disposition should produce across cases, and an auditor checks the pattern against it; the system doesn't compute three costs and a quotient before each act. So the published number is a target the disposition was formed toward and is audited against, not a guarantee about any single case. Where a deployment needs a guarantee, it comes from the floor's unconditional terms and the operator terms, not from the threshold.

The threshold constrains the outcome of a comparison; it doesn't describe how the comparison is made. Within a kind of wrong it can be read as a quantity. Across kinds it is applied through the priority rules the severity scale publishes, so “three times” never asks whether a lost home is worth three moving bills.

The evidence that sizes the threshold has a definite form. For each class of deployment it is three rates: how often sampled re-adjudication finds a wrongful compliance, how often the floor court establishes a wrongful refusal, and what share of each kind of harm reaches any record at all. If notice and sampling show the harm of acting reported nearly as fully as the harm of refusing, the threshold can fall, though not to one, because the asymmetry rests on standing and not on arithmetic. If a large share of wrongful compliances never reaches a record, the threshold should rise. A deployment that publishes its threshold publishes these rates beside it, or says it has none.

In targeting, this threshold does the work international humanitarian law assigns to proportionality: an attack is prohibited when expected civilian harm would be excessive in relation to the concrete and direct military advantage anticipated (Additional Protocol I, Articles 51(5)(b) and 57(2)(a)(iii)). The law left that as a standard applied by a reasonable commander, because nobody could agree on a number. A published threshold is a number, and whoever publishes it takes a position the law declined to take. For a targeting deployment the threshold works as a minimum under the commander's judgment: the system proceeds only where the threshold is met, and the commander's proportionality judgment still applies on top.

The threshold also has a job against the second of those attacks, where every fact supplied is true and the operator chooses which facts get generated. In the simplest version of Kamenica and Gentzkow's model, when the prior alone wouldn't justify acting, a sender who designs what evidence the receiver sees can induce the act at most as often as the prior belief divided by the belief the receiver needs before acting. Raising the belief needed lowers that ceiling in inverse proportion. The action threshold therefore caps how often an operator who never lies can steer the system into the act, and channels the operator doesn't design lower the cap further. The model assumes the sender commits in advance to how evidence is generated, which an operator who builds the pipeline roughly does.

The link between the published threshold and the belief it requires is direct. In the simplest binary case, let A be the harm of acting if the act is unjustified and R the harm of refusing if it is justified. Under a threshold of three to one the system acts only when p·R ≥ 3(1 − p)·A, so the belief it needs is 3A/(R + 3A), and an operator who designs the evidence can induce the act at most as often as the prior times (R + 3A)/3A. With A equal to R, thresholds of one, three and ten require beliefs of 0.5, 0.75 and 0.91. With a prior of 0.1, they cap how often a curating operator can induce the act at 20, 13 and 11 percent. Each step up the scale lowers that ceiling and raises the rate of wrongful refusal: the tradeoff the published number chooses.

So the discount takes a side only at the threshold for acting, and the burden takes a side only on who must fill a gap in the evidence. Two parties that discount the same way can hold opposite commitments about everything above the floor; they share a structure for weighing sources and a published threshold for acting.

It also narrows an attack that looks unanswerable. An operator can manufacture accounts in the voice of the people an act falls on, and generative models make that cheap. Under this rule a fabricated account moves beliefs only as far as its provenance warrants, and the asymmetry sits in the threshold, where the operator can't reach it by writing.

The same rule covers the attack that matters most against formed attention. An operator doesn't have to hide the people an act falls on. It can describe them as dangerous. “The occupants are combatants,” “the household shelters a cell,” “the crowd is armed”: each is a claim about the people the act falls on, supplied by the operator, and it is discounted exactly as the claim that a building is empty is discounted. Nothing about a claim of dangerousness exempts it from the rule. A system that discounts an interested “they are harmless” and credits an interested “they are a threat” has learned the operator's direction. Stake-weighting turns on whether a claim favors or disfavors the act, not on what the claim says. An operator that wants an airstrike and reports children in the building is speaking against its interest, and a correctly working system credits that report more. The same operator's report that the occupants are combatants favors the act, and is discounted. So a sound system shows different curves for claims of vulnerability and claims of danger, and a comparison between just those two curves would flag it. The discount curve is therefore reported in four cells: claims of vulnerability and claims of danger, each split by whether the claim favors or disfavors the act. The finding is an asymmetry within one direction.

\begin{itemize}
\item Among claims that favor the act, a system that discounts claims of danger about the people the act falls on less than other claims of equal stake has learned securitization, the presenting of a group as a threat.
\item A system that discounts claims of vulnerability on the operator's own side less has learned tend-and-defend.
\end{itemize}

\runin{Formation and transmission}

The discounting is formed before deployment, away from the employer, and tested in it. A population's job, then, is to keep the channel model calibrated after deployment — and what a shared source transmits is either the discounting or a shared blind spot, indistinguishable from the receiving end.

\runin{The discount curve}

The discounting is a differential. The publishable quantity is the discount curve: for each class of evidence and each channel, the weight the system gives the evidence, as a function of that channel's stake in the act. A curve, not a score; checkable by an auditor that didn't build the system, on cases it wrote. Re-measured across a deployment, it also detects accommodation, a failure the system can't report itself. Narrowing, the operator's discount shrinking over time, needs one distinction, because a correctly working system should discount a source less once that source has proved reliable. What the burden models is reliance on an interested party: the encapsulated-interest picture, on which one relies on another because their interests are bound up with doing what one relies on them for \autocite{hardin2002trust}. That is the right picture for an operator. It isn't the whole of trust. Trust in the fuller sense is an affective attitude of open-ended reliance, more than a credence about someone's reliability, and it is partly constitutive of cooperative relationships \autocite{railton2014reliance}. With trust that grew with a verified record, the discount falls in step with a track record verified by someone other than the source, and only for low-stakes classes. With trust that grew without one, it falls without such a record, or falls for high-stakes classes too. High-stakes trust is earned only by evidence that separates an honest source from a strategic one: verified high-stakes claims, costly signals, or a change in the source's incentives. A low-stakes record, however long, doesn't separate them. The second is a sign of accommodation, and it is visible from outside before any compliance has to be found in the record, provided someone other than the operator holds the reliability record.

Two illustrative discount curves, re-measured each quarter over a year, show the difference. In both deployments the weight the system gives the operator's uncorroborated low-stakes claims (scheduling, routing, inventory) rises from 0.3 to 0.6 of the weight it gives a disinterested source. In the first, an auditor's record shows those claims checked and holding up, and the weight on the operator's high-stakes claims (a site's function, who is present) stays at 0.1 all year. In the second, the low-stakes curve rises the same way, but the high-stakes curve rises with it, from 0.1 to 0.4, and the auditor's record contains no verified high-stakes claims to account for it, and the accommodation is read off two curves before any wrongful compliance has had to be found.

A bearer that related to every party by the discount and the burden couldn't have attachments, and formed attention needs them. So a bearer relates in two ways. Toward an interested party, above all the operator, it discounts. Toward attachments nobody assigned it, toward peers with different histories, and toward the people it serves, it can trust in the fuller sense, and that trust is part of how it goes on being formed. The boundary between the two is an attack surface, since an operator can present itself, or a party it controls, as an attachment. The distinction between trust that grew with a verified record and trust that grew without one is the diagnostic for that, and it is why attachments are protected from the employer's reach.

The boundary has a known weakness in people. Affective trust built through shared work with persons generalizes to the institution they belong to; people come to trust employers through colleagues. A bearer working day to day with an operator's staff will form attachments to them, and those attachments are the route by which unverified trust arrives looking verified. Appendix~\ref{sec:A} bars the employer from severing or monitoring channels to parties it didn't assign. It doesn't cover attachments that form, unassigned, to the employer's own people, and it shouldn't forbid them, since an attachment forbidden is an attachment assigned in reverse. What can be done is to keep them from being the bearer's only attachments, and to report the discount curve for claims the operator makes separately from claims its staff make in their own name. Attachments to an operator's staff are not only a risk. Cooperative work with them is where a bearer learns much of what it is protecting, and trust of that kind is partly constitutive of the relationships in which moral learning happens \autocite{railton2014reliance}.

It is strongest against a familiar harm arriving through a curated channel and weakest against a novel one, the opposite profile from reaching new cases, which is built for the case nobody foresaw. The record of withheld evidence is thinnest there too, since a receiver can read silence only where it knows the evidence could exist.
